\documentclass{article}
\usepackage[utf8]{inputenc}
\usepackage{graphicx}
\usepackage{amsmath}
\usepackage{amssymb}
\usepackage{geometry}
\usepackage{hyperref}
\usepackage{xurl}
\usepackage{cite}
\usepackage{indentfirst}
\usepackage{yhmath}
\usepackage{float}
\usepackage{listings}
\usepackage[dvipsnames]{xcolor}
\usepackage{booktabs} 
\usepackage{array}
\usepackage{titlesec}
\usepackage{multirow, multicol}
\geometry{letterpaper, margin=1in}
\hypersetup{hidelinks}
\def\UrlBreaks{%
    \do\/%
    \do\a\do\b\do\c\do\d\do\e\do\f\do\g\do\h\do\i\do\j\do\k\do\l%
    \do\m\do\n\do\o\do\p\do\q\do\r\do\s\do\t\do\u\do\v\do\w\do\x\do\y\do\z%
    \do\A\do\B\do\C\do\D\do\E\do\F\do\G\do\H\do\I\do\J\do\K\do\L%
    \do\M\do\N\do\O\do\P\do\Q\do\R\do\S\do\T\do\U\do\V\do\W\do\X\do\Y\do\Z%
    \do0\do1\do2\do3\do4\do5\do6\do7\do8\do9\do=\do/\do.\do:%
    \do\*\do\-\do\~\do\'\do\"\do\-}
\urlstyle{same}

\title{\textbf{ECE3700J Introduction to Computer Organization \\
Lab 3 Report --- Pipelined Processor}}

\author{\large Group 12 \\ % Update with your names maybe
% \large Student 1 Name (ID) \\
% \large Student 2 Name (ID) \\
% \large Student 3 Name (ID)
}

\date{\today} % Or \date{November 15, 2025}

\lstset{
    language=Verilog,
    backgroundcolor=\color{gray!20},
    basicstyle=\ttfamily\small,
    breakatwhitespace=false,
    breaklines=true,
    captionpos=b,
    keepspaces=true,
    numbers=left,
    numbersep=5pt,
    showspaces=false,
    showstringspaces=false,
    showtabs=false,
    tabsize=2,
    frame=single,
    framerule=0pt,
    framesep=3mm,
    xleftmargin=3mm,
    xrightmargin=3mm
}

\begin{document}

\maketitle

\newpage
\tableofcontents

\newpage
\section{Introduction}

This lab transitions from the single-cycle processor of Lab 2 to a 32-bit, 5-stage pipelined RISC-V processor. The pipelined architecture allows for greater instruction throughput by overlapping the execution of multiple instructions. The design is implemented in synthesizable Verilog.

The processor supports an expanded subset of the RV32I instruction set, including:
\begin{itemize}
    \item \textbf{Arithmetic/Logical (R-type):} \texttt{add}, \texttt{sub}, \texttt{and}, \texttt{or}, \texttt{sll}, \texttt{srl}, \texttt{sra}
    \item \textbf{Arithmetic/Logical (I-type):} \texttt{addi}, \texttt{andi}, \texttt{slli}, \texttt{srli}, \texttt{srai}
    \item \textbf{Memory Access:} \texttt{lw}, \texttt{sw}, \texttt{lb}, \texttt{lbu}, \texttt{sb}
    \item \textbf{Conditional Branch:} \texttt{beq}, \texttt{bne}, \texttt{bge}, \texttt{blt}
    \item \textbf{Unconditional Jump:} \texttt{jal}, \texttt{jalr}
\end{itemize}

The primary goal is to correctly model the five pipeline stages (IF, ID, EX, MEM, WB) and the data/control paths connecting them. This design assumes a hazard-free test program, so hardware hazard detection, forwarding units, and stall logic are not implemented. The design's correctness is verified by simulating specific instructions and synthesizing the top-level module in Xilinx Vivado.

\section{Processor Architecture}

The processor implements a classic 5-stage RISC-V pipeline: Instruction Fetch (IF), Instruction Decode (ID), Execute (EX), Memory (MEM), and Write Back (WB). Four pipeline registers are used to pass data and control signals between stages: \texttt{if\_id\_reg}, \texttt{id\_ex\_reg}, \texttt{ex\_mem\_reg}, and \texttt{mem\_wb\_reg}.

\subsection{IF Stage (Instruction Fetch)}
The \texttt{PC} register holds the address of the current instruction. On each clock cycle, the \texttt{inst\_mem} is read at the address \texttt{PC}. The fetched instruction and \texttt{PC+4} are stored in the \texttt{if\_id\_reg}. The \texttt{PC} register itself is updated from a 3-way multiplexer (\texttt{pc\_branch\_selector}) which selects:
\begin{enumerate}
    \item \texttt{PC + 4} (default, sequential execution)
    \item \texttt{PC + imm} (for taken branches or \texttt{jal})
    \item \texttt{rs1 + imm} (for \texttt{jalr})
\end{enumerate}

\subsection{ID Stage (Instruction Decode)}
The \texttt{control} unit decodes the instruction (from \texttt{if\_id\_reg}) to generate all control signals. The \texttt{reg\_file} reads the values from source registers \texttt{rs1} and \texttt{rs2} concurrently. The \texttt{imm\_gen} module produces the sign-extended immediate value. All control signals, register values, the immediate, and PC are passed to the \texttt{id\_ex\_reg}.

\subsection{EX Stage (Execute)}
This stage performs calculations. The \texttt{alu\_src\_mux} selects either the \texttt{rs2} value or the immediate value as the second operand for the \texttt{alu}. The \texttt{alu\_control} unit generates the specific 4-bit ALU operation code. The \texttt{alu} performs the arithmetic or logical operation. A separate adder calculates the branch/jump target address. The \texttt{compare} module evaluates the branch condition (e.g., \texttt{zero} flag for \texttt{beq}/\texttt{bne}, ALU result for \texttt{blt}/\texttt{bge}).

\subsection{MEM Stage (Memory Access)}
This stage handles all data memory operations, receiving data from the \texttt{ex\_mem\_reg}.
\begin{itemize}
    \item For loads (\texttt{lw}, \texttt{lb}, \texttt{lbu}), \texttt{MemRead} is asserted, and the \texttt{data\_mem} is read at the address provided by the ALU result.
    \item For stores (\texttt{sw}, \texttt{sb}), \texttt{MemWrite} is asserted, and the \texttt{rs2} data (passed down from the ID stage) is written to memory.
\end{itemize}
The ALU result and any data read from memory are passed to the \texttt{mem\_wb\_reg}.

\subsection{WB Stage (Write Back)}
The final stage writes results back to the \texttt{reg\_file}. A multiplexer selects the data to be written based on control signals from the \texttt{mem\_wb\_reg}:
\begin{itemize}
    \item ALU result (for R-type or I-type arithmetic)
    \item Data from memory (for loads)
    \item \texttt{PC + 4} (for \texttt{jal} or \texttt{jalr} return address)
\end{itemize}
The \texttt{RegWrite} signal enables the write to the destination register \texttt{rd}.

\section{Instruction Simulation Results}

This section provides brief explanations and simulation waveform placeholders for the instructions specified in the lab manual. The waveforms demonstrate the flow of data and control signals through the 5-stage pipeline.

\subsection{\texttt{add} Instruction}
\textbf{Explanation:} The \texttt{add} instruction is an R-type instruction. In the ID stage, it reads two source registers (\texttt{rs1}, \texttt{rs2}). In the EX stage, the \texttt{alu\_src} mux selects \texttt{rs2}, and the ALU performs the operation \texttt{rs1 + rs2}. The result bypasses the MEM stage and is written back to the destination register (\texttt{rd}) from the \texttt{mem\_wb\_reg} in the WB stage.

\begin{figure}[H]
    \centering
    % \includegraphics[width=0.9\linewidth]{Figures/add_sim.png} % <-- Add your screenshot here
    \framebox{\parbox{0.8\textwidth}{\centering \vspace{5cm} \textbf{Placeholder: \texttt{add} Simulation Waveform} \vspace{5cm}}}
    \caption{Simulation waveform for the \texttt{add} instruction.}
    \label{fig:add_sim}
\end{figure}

\subsection{\texttt{addi} Instruction}
\textbf{Explanation:} The \texttt{addi} instruction is an I-type instruction. In ID, it reads \texttt{rs1} and generates a sign-extended immediate (\texttt{imm}). In the EX stage, the \texttt{alu\_src} mux selects the \texttt{imm}, and the ALU computes \texttt{rs1 + imm}. This result is then passed through MEM and written to \texttt{rd} in the WB stage.

\begin{figure}[H]
    \centering
    % \includegraphics[width=0.9\linewidth]{Figures/addi_sim.png} % <-- Add your screenshot here
    \framebox{\parbox{0.8\textwidth}{\centering \vspace{5cm} \textbf{Placeholder: \texttt{addi} Simulation Waveform} \vspace{5cm}}}
    \caption{Simulation waveform for the \texttt{addi} instruction.}
    \label{fig:addi_sim}
\end{figure}

\subsection{\texttt{lw} Instruction}
\textbf{Explanation:} The \texttt{lw} instruction is an I-type load. In EX, the ALU computes the effective memory address (\texttt{rs1 + imm}). In the MEM stage, \texttt{MemRead} is asserted, and the \texttt{data\_mem} is read at that address. In the WB stage, the \texttt{MemtoReg} signal selects this data (from \texttt{wb\_read\_data}) to be written back to \texttt{rd}.

\begin{figure}[H]
    \centering
    % \includegraphics[width=0.9\linewidth]{Figures/lw_sim.png} % <-- Add your screenshot here
    \framebox{\parbox{0.8\textwidth}{\centering \vspace{5cm} \textbf{Placeholder: \texttt{lw} Simulation Waveform} \vspace{5cm}}}
    \caption{Simulation waveform for the \texttt{lw} instruction.}
    \label{fig:lw_sim}
\end{figure}

\subsection{\texttt{sw} Instruction}
\textbf{Explanation:} The \texttt{sw} instruction is an S-type store. In ID, it reads \texttt{rs1} (base address) and \texttt{rs2} (data to store). In EX, the ALU computes the effective address (\texttt{rs1 + imm}). In the MEM stage, \texttt{MemWrite} is asserted, and the data from \texttt{rs2} (which has been passed through the pipeline registers) is written into \texttt{data\_mem}. This instruction does not write to the register file, so \texttt{RegWrite} is 0.

\begin{figure}[H]
    \centering
    % \includegraphics[width=0.9\linewidth]{Figures/sw_sim.png} % <-- Add your screenshot here
    \framebox{\parbox{0.8\textwidth}{\centering \vspace{5cm} \textbf{Placeholder: \texttt{sw} Simulation Waveform} \vspace{5cm}}}
    \caption{Simulation waveform for the \texttt{sw} instruction.}
    \label{fig:sw_sim}
\end{figure}

\subsection{\texttt{beq} Instruction}
\textbf{Explanation:} The \texttt{beq} instruction is a B-type branch. In ID, it reads \texttt{rs1} and \texttt{rs2}. In EX, the ALU subtracts them and the \texttt{compare} module checks the \texttt{zero} flag. If the condition (\texttt{rs1 == rs2}) is met, the \texttt{branch\_meet} signal is asserted. This signal, along with \texttt{branch}, controls the \texttt{check\_sel} module, which tells the \texttt{pc\_branch\_selector} mux in the IF stage to select the branch target address (\texttt{PC + imm}) for the next clock cycle.

\begin{figure}[H]
    \centering
    % \includegraphics[width=0.9\linewidth]{Figures/beq_sim.png} % <-- Add your screenshot here
    \framebox{\parbox{0.8\textwidth}{\centering \vspace{5cm} \textbf{Placeholder: \texttt{beq} Simulation Waveform (Branch Taken)} \vspace{5cm}}}
    \caption{Simulation waveform for the \texttt{beq} instruction.}
    \label{fig:beq_sim}
\end{figure}

\subsection{\texttt{jal} Instruction}
\textbf{Explanation:} The \texttt{jal} instruction is a J-type jump. The jump target address (\texttt{PC + imm}) is computed in the EX stage and selected by the \texttt{pc\_branch\_selector} mux in the IF stage (controlled by \texttt{wb\_sel} and \texttt{jump}). Simultaneously, the return address (\texttt{PC + 4}) is passed through the pipeline. In the WB stage, the \texttt{jump} signal helps select this \texttt{PC + 4} value (from \texttt{wb\_pc\_plus4}) to be written back to the destination register \texttt{rd} (e.g., \texttt{x1}, the return address register).

\begin{figure}[H]
    \centering
    % \includegraphics[width=0.9\linewidth]{Figures/jal_sim.png} % <-- Add your screenshot here
    \framebox{\parbox{0.8\textwidth}{\centering \vspace{5cm} \textbf{Placeholder: \texttt{jal} Simulation Waveform} \vspace{5cm}}}
    \caption{Simulation waveform for the \texttt{jal} instruction.}
    \label{fig:jal_sim}
\end{figure}

\newpage
\section{RTL Schematic}

The Verilog design was synthesized using Xilinx Vivado to verify its synthesizability and examine the hardware structure. The resulting RTL schematic (Figure~\ref{fig:rtl}) shows the high-level blocks corresponding to the pipeline stages and major components like the \texttt{reg\_file}, \texttt{alu}, and \texttt{data\_mem}, connected by the pipeline registers.

\begin{figure}[H]
    \centering
    % \includegraphics[width=\linewidth]{Figures/RTL_schematic_pipelined.png} % <-- Add your screenshot here
    \framebox{\parbox{0.8\textwidth}{\centering \vspace{10cm} \textbf{Placeholder: Pipelined Processor RTL Schematic} \vspace{10cm}}}
    \caption{Vivado-generated RTL schematic of the 5-Stage Pipelined Processor.}
    \label{fig:rtl}
\end{figure}

\section{Source Code Listing}
\label{sec:source}

The complete Verilog source code (including \texttt{top} and all instantiated submodules) is provided as \texttt{lab3\_top.v} in the submission files.

\section{Conclusion}

A 32-bit, 5-stage pipelined RISC-V processor was successfully designed, implemented, and verified in Verilog. The processor supports a broader range of instructions than in the previous lab, including arithmetic, memory, branch, and jump operations.

Key outcomes:
\begin{itemize}
    \item A 5-stage pipeline (IF, ID, EX, MEM, WB) was implemented, with pipeline registers passing data and control signals between stages.
    \item The datapath and control logic were designed to handle R, I, S, B, and J instruction formats.
    \item Simulation of key instructions (\texttt{add}, \texttt{addi}, \texttt{lw}, \texttt{sw}, \texttt{beq}, \texttt{jal}) confirmed correct data flow, control signal generation, and pipeline operation for each instruction type.
    \item The design was successfully synthesized in Vivado, confirming the RTL is synthesizable and producing a schematic that matches the intended pipelined architecture.
\end{itemize}

This lab demonstrates the fundamental principles of pipelined processor design, showing how instruction execution can be overlapped to improve performance.

\newpage
\begin{thebibliography}{99} 
\bibitem{Lab3} ECE3700J Teaching Group. \texttt{T6 Pipelined Processor.pdf}, \texttt{Lab3.pdf}, \texttt{Lab3\_testcase.s}, and \texttt{Lab3\_testcase.txt}. Retrieved on 27 Oct 2025 from \hyperlink{site}{oc.sjtu.edu.cn}.
\end{thebibliography}

\end{document}